\documentclass[11pt]{amsart}

\usepackage{amsmath,amssymb,amsthm}
\usepackage[margin=1.2in]{geometry}
\usepackage{microtype}
\usepackage{tikz}
\usepackage[colorlinks=true,linkcolor=blue!50!black,citecolor=blue!50!black,urlcolor=blue!50!black]{hyperref}

\newtheorem{theorem}{Theorem}
\newtheorem{lemma}[theorem]{Lemma}
\newtheorem{proposition}[theorem]{Proposition}
\theoremstyle{definition}
\newtheorem{definition}[theorem]{Definition}
\theoremstyle{remark}
\newtheorem{remark}[theorem]{Remark}

\title{The square achievement game on $n \times n$ boards}
\author{Sean Reid}
\date{August 2026}

\begin{document}

\begin{abstract}
Two players alternately claim cells of an $n \times n$ grid; the first
to hold all four corners of an axis-aligned square wins. We prove the
first player has a winning strategy if and only if $n \ge 5$. The
strategy is explicit: two reflexes over a finite table of forcing plans
confined to a corner region of the board, so a single certificate
covers every sufficiently large board. All results, including the
exhaustive analyses of small boards, are formalized and machine-checked
in Lean~4.
\end{abstract}

\maketitle

\section{Introduction}

% History and positioning: Harary's achievement games, positional game
% theory (Beck), Maker--Maker conventions and why draws matter, prior
% work on polyomino and square achievement, statement of results.

The square achievement game is played on the cells of an $n \times n$
grid. Two players alternately claim empty cells, the first player (P1)
moving first, and a player wins on completing the four corners
\[
  \{(r,c),\ (r,c+s),\ (r+s,c),\ (r+s,c+s)\}, \qquad s \ge 1,
\]
of an axis-aligned square. If the board fills with no completed square
the game is drawn.

\begin{theorem}\label{thm:main}
P1 has a winning strategy in the square achievement game on the
$n \times n$ board if and only if $n \ge 5$.
\end{theorem}

\section{The game and conventions}

% Formal definitions matching the Lean development: cells in the
% quarter-plane, boards as finite subgrids, winning sets, the game tree,
% draws counting against the first player for the purposes of the
% threshold statement.

\section{Small boards}\label{sec:small}

% Exhaustive refutation for n <= 4 with the pruned dual search; the
% strategy-stealing argument showing P2 never wins, hence draws.

\section{The winning strategy}\label{sec:strategy}

% The reflex-plus-table architecture; forcing plans; the corner-box
% confinement; the n = 5 special case and why pure forcing fails there;
% the alignment classes for far replies.

\section{Search and extraction}\label{sec:pipeline}

% How the tables were found: the forcing-script search, the exact
% solver for n = 5, extraction of strategy trees, and the two
% independent verifiers (Python and compiled Lean) used during
% development.

\section{The saturation argument}\label{sec:saturation}

% The subsumption theorem: activity is confined, far stones interact
% through finitely many alignment classes, and one verified bound
% covers all larger boards.

\section{Formalization}\label{sec:lean}

% The Lean architecture: the game as an inductive winning predicate,
% executable checkers with machine-checked soundness in both
% directions, certificates discharged by native_decide, and the axiom
% footprint.

\section{Reproducibility}

% Repository layout, how to regenerate tables, re-run every
% verification, and rebuild this paper.

\begin{thebibliography}{9}

\bibitem{harary}
F.~Harary,
\emph{Achievement and avoidance games for graphs},
Ann. Discrete Math. \textbf{13} (1982), 111--119.

\bibitem{beck}
J.~Beck,
\emph{Combinatorial Games: Tic-Tac-Toe Theory},
Encyclopedia of Mathematics and its Applications 114,
Cambridge University Press, 2008.

\bibitem{lean4}
L.~de~Moura and S.~Ullrich,
\emph{The Lean 4 theorem prover and programming language},
CADE 28, 2021.

\bibitem{mathlib}
The mathlib Community,
\emph{The Lean mathematical library},
CPP 2020.

\end{thebibliography}

\end{document}
